\documentclass[11pt]{article}
\usepackage{amsmath}
\usepackage{amssymb}
\usepackage{amsthm}
\usepackage{enumitem}

\begin{document}
\newcommand{\seq}[2][n]{\ensuremath{\left\{ #2 _{#1} \right\}}}
\section{Thorem 10.2}
    Let $(X,d)$ be a compact metric space. 
    Then, every sequence in $(X,d)$ has a convergent subsequence.
    \begin{proof}
        Let $(X,d)$ be compact and let $\{x_n\}$ be any sequence in $X$.

        \textbf{Case 1:} \textit{$\{x_n\}$ has only finitely many values.}

            Then, by the pigeonhole principle, there exists $x \in X$ such that $x_n = x$ infinitely often.
            This gives a subsequence $\{x_{n_k}\}$ that converges to $x$.

        \textbf{Case 2:} \textit{$\{x_n\}$ has infinitely many values.}

            By Lemma 10.4, $\{x_n\}$ has a limit point $x \in X$.
            Since $x$ is a limit point in $\{x_n\}$, 
            \begin{itemize}
                \item   There exists $n_1 \in \mathbb{N}$ such that $d(x_{n_1},x) < 1$
                \item   There exists $n_2 \in \mathbb{N}$ such that $d(x_{n_2},x) < \frac{1}{2}$ where $n_2 > n_1$ and $x_{n_1} \neq x$.
                \item   For all $k \in \mathbb{N} \backslash \{1\}$, there exists $n_k \in \mathbb{N}$ such that $d(x_{n_k},x) < \frac{1}{k}$ where $n_k > n_{k-1}$ using $x_{n_k} \neq x$.
            \end{itemize}
            By construction, $\{x_{n_k}\}$ is a subsequence of $\{x_n\}$ and $x_{n_k}$ converges to $x$.
    \end{proof}

\section{Lemma 10.4}
    Let $(X,d)$ be compact and let $E \subset X$ be infinite.
    Then $E$ has a limit point in $X$.
    \begin{proof}[Proof by Contradiction]
        Suppose for contradiction that $E$ has no limit point in $X$.
        Then, for all $x \in X$, $x$ is not a limit point of $E$.
        This means that we can find $r_x > 0$ such that $B_{r_x}(x)$ contains no points of $E$ besides maybe $x$.
        Therefore, each $B_{r_x}(x)$ contains at most one one point of $E$.
        Since this is for all $x \in X$, $\left\{ B_{r_x}(x) \right\}_{x \in X}$ is an open cover of $X$.
        Since $X$ is compact, $\left\{ B_{r_x}(x) \right\}_{x \in X}$ has a finite subcover containing at most one element of $X$ each and thereby at most one element of $E$ each.
        Therefore, $E$ is finite.
        This contradicts that the cardinality of $E$ is infinite.
    \end{proof}

\section{Theorem 10.8}
    \begin{enumerate}[label=(\alph*)]
        \item   Every compact metric space is complete
        \item   $\mathbb{R}$ is complete
    \end{enumerate}
    \begin{proof}[Proof (a)]
        A complete metric has every Cauchy sequence converge.
        Let $(X,d)$ be compact and let \seq{x} be cauchy.
        By Theorem 10.2, \seq{x} has a convergent subsequence.
        By Lemma 10.9, \seq{x} converges.
        Therefore, $X$ is complete.
    \end{proof}
    \begin{proof}[Proof (b)]
        In $\mathbb{R}$, every convergent sequence \seq{x} is bounded, meaning it is contained in some $[a,b]$.
        $[a,b]$ is compact and \seq{x} is bounded, so it has a convergent subsequence and converges, so $\mathbb{R}$ is complete.
    \end{proof}

\section{Lemma 10.9}
    Let \seq{x} be cauchy.
    If \seq{x} has a convergent subsequence, it converges.
    \begin{proof}
        Suppose that \seq[n_\ell]{x} is a subsequence of \seq{x} that converges to $x$.
        Therefore, given $\varepsilon > 0$, there exists $K \in \mathbb{N}$ such that if $k \geq K$, $d(x_{n_k}, x) < \frac{\varepsilon}{2}$.
        Similarly, since \seq{x} is cauchy, there exists some $N \in \mathbb{N}$ such that if $m \geq n \geq \mathbb{N}$, $d(x_m, x_n) < \frac{varepsilon}{2}$.
        Pick $K$ big enough that $n_K \geq N$.
        This means that if $m \geq n_K$, then $\frac{\varepsilon}{2} + \frac{\varepsilon}{2} = \varepsilon > d(x_{n_k}, x) + d(x_m, x_{n_k}) \geq d(x_m, x)$.
        Therefore, \seq{x} converges to $x$.
    \end{proof}
\end{document}